% =============================================================================
% 05_report_quality.tex — Problem 5: Report Quality and Reproducibility
% =============================================================================
\section{Report Quality and Reproducibility}
\label{sec:report-quality}

% ---- 5.1 Reproduction Record ------------------------------------------------
\subsection{Reproduction Record}
\label{sec:reproduction}

The following step-by-step instructions reproduce all results, figures, and
tables in this report from scratch.

\paragraph{Prerequisites.}
\begin{itemize}
    \item R $\geq$ 4.3.0 with packages: \texttt{glmnet}, \texttt{ggplot2},
          \texttt{dplyr}, \texttt{corrplot}, \texttt{caret}, \texttt{knitr},
          \texttt{xtable}.
    \item \LaTeX\ distribution (e.g.\ TeX Live 2024 or MiKTeX) with
          \texttt{latexmk}.
    \item The dataset \texttt{fat.csv} placed in
          \texttt{project\_quiz1\_latex/data/}.
\end{itemize}

\paragraph{Steps.}

\begin{lstlisting}[language=bash, caption={Reproduction commands.}]
# 1. Clone or download the project
cd D:/Stats2/.

# 2. Run the R analysis pipeline
make all
# This generates all figures in output/figures/
# and all tables in output/tables/

# 3. Compile the LaTeX report
cd report
latexmk -pdf main.tex

# 4. Clean auxiliary files (optional)
latexmk -c
\end{lstlisting}

% TODO: Add any additional steps or notes about the environment setup,
%       e.g., specific R package versions, OS-specific instructions.

% ---- 5.2 Recommendation and Limitations ------------------------------------
\subsection{Recommendation and Limitations}
\label{sec:recommendation}

% TODO: Write a final recommendation addressing:
%   1. Which model should be used for predicting brozek body-fat percentage?
%   2. Is the goal prediction accuracy or interpretability? How does this
%      affect the recommendation?
%   3. What are the key limitations of this analysis?
%      (e.g., small sample size, linearity assumption, limited predictor set,
%       potential measurement error, generalisability to other populations)

\paragraph{Recommendation.}
\textit{TODO: State your final model recommendation with supporting evidence
from the cross-validation and holdout results.}

\paragraph{Prediction vs.\ Interpretation.}
\textit{TODO: Discuss the trade-off. Ridge retains all predictors (better for
interpretation of relative importance), while Lasso provides a sparse model
(easier to deploy and communicate). Elastic Net offers a middle ground.}

\paragraph{Limitations.}
\begin{enumerate}
    \item \textit{TODO: Limitation 1 (e.g., sample size and generalisability).}
    \item \textit{TODO: Limitation 2 (e.g., linearity assumption).}
    \item \textit{TODO: Limitation 3 (e.g., predictor measurement error).}
\end{enumerate}

% ---- 5.3 AI Verification Record --------------------------------------------
\subsection{AI Verification Record}
\label{sec:ai-verification}

All AI-assisted outputs were independently verified by team members.
Table~\ref{tab:ai-verification} summarises the verification process.

\begin{table}[H]
\centering
\caption{AI verification record.}
\label{tab:ai-verification}
\begin{tabularx}{\textwidth}{l X X l}
\toprule
\textbf{AI-Assisted Item} &
\textbf{How Checked} &
\textbf{Evidence} &
\textbf{Modification} \\
\midrule
% TODO: Fill in each row. Examples:
%
% Ridge derivation &
% Verified each step by hand &
% Handwritten notes (scanned) &
% None \\[4pt]
%
% R code for CV &
% Compared output with manual calculation on subset &
% Console output matches &
% Fixed seed for reproducibility \\[4pt]
%
\textit{TODO: Item 1} & \textit{TODO} & \textit{TODO} & \textit{TODO} \\[4pt]
\textit{TODO: Item 2} & \textit{TODO} & \textit{TODO} & \textit{TODO} \\[4pt]
\textit{TODO: Item 3} & \textit{TODO} & \textit{TODO} & \textit{TODO} \\
\bottomrule
\end{tabularx}
\end{table}
